\section{代码审查的价值与方法}

\subsection{为什么代码审查是最高杠杆的工程活动}

代码审查（Code Review）是提升个人工程能力和团队代码质量最具杠杆效应的活动之一。数据佐证：

\begin{importantbox}{代码审查的量化价值}
\begin{itemize}
    \item 代码审查的错误检出率为 55--60\%，而各种测试模式的检出率仅为 25--45\%
    \item 无代码审查 vs 有代码审查的程序错误率：4.5 $\rightarrow$ 0.82 errors/100 lines
    \item AT\&T 的研究显示：代码审查带来 14\% 的生产力提升和 90\% 的缺陷减少
\end{itemize}
\textit{来源：Coding Horror}
\end{importantbox}

\subsection{代码审查应该关注什么}

\subsubsection{逻辑与正确性错误}

代码是否按预期工作？是否有边界条件处理不当、off-by-one 错误等？

\subsubsection{可读性与可维护性}

代码是否易于理解和后续修改？命名是否清晰？函数职责是否单一？

\subsubsection{性能}

是否存在不必要的数据库查询、N+1 问题、不合理的时间/空间复杂度？

\subsubsection{安全}

是否存在 SQL injection、XSS、未验证的输入、敏感数据泄露等安全隐患？

\subsubsection{最佳实践}

是否遵循项目的编码规范、设计模式和架构约定？

\subsection{什么是好的代码审查}

\begin{warningbox}{``This won't work'' 不是好的审查意见}
简单的否定不提供任何教学价值。好的审查意见应该是建设性的、具体的，并邀请对话。例如：

``I see your new method matches the existing style in this file, taking {[}X{]} parameters. Having that many parameters hurts readability and implies the function is doing too much. What do you think about refactoring this method and the existing ones in a later pull request to reduce how many parameters they take?''
\end{warningbox}

\subsection{本章小结}

代码审查是软件质量的核心保障，其错误检出率远超自动化测试。好的审查意见应该是具体的、建设性的，并为被审查者提供学习机会。

